% Checked with grammarly
\section{Tech stack}

For the development, Visual Studio Code was used due to the rich extensions marketplace and easy integration of C, C++ and Java programs (https://code.visualstudio.com/). Code formatting is handled by ClangFormat with a custom .clang-format configuration file to enforce an consistent style guide for improving readability and for maintaining a cleaner code base(https://clang.llvm.org/docs/ClangFormat.html).The main language which was used is C++17 and c11 due to the reasons explained in 3.1 Used Programming language.  For building and linking of the .cpp and .hpp source file the GNU C++ compiler (GCC 13.2.0) was used (https://gcc.gnu.org/). The build process is automated via a Makefile, which can be executed using \texttt{mingw32-make} on windows, \texttt{add here for linux} on Linux and \texttt{add here for mac} on Mac OS. Then to run the executable programm ./anomaly\_detection or just right away building and running by adding "run" to the make script shown before.
For the version control Git 2.53.0 is used and whole code is documented on GitHub with an clear introductions on how to use the code :\url{https://github.com/FabianSanterTUGraz/Real-time-Anomaly-Detection-in-Time-Series-Data}.While the \texttt{matplotlib.pyplot} library was utilized for initial data visualization, plotting functionality serves solely as a visual diagnostic tool for inspecting and comparing results, and is therefore excluded from the core algorithm specifications.